\section{Формализация объектов сцены}

Размеры задаются в условных единицах (у.е.), ось $z$ направлена вертикально вверх, сцена занимает по горизонтали квадрат $[0;\,100] \times [0;\,100]$~у.е. Сцена содержит следующие объекты.
\begin{enumerate}
	\item Ландшафт "--- заснеженная поверхность высотой $H(x, y) \in [0;\,10]$~у.е., заданная картой высот (полутоновым изображением).
	\item Статические объекты типов <<дом>>, <<ель>> и <<автомобиль>>, составленные из стандартных тел (параллелепипедов, пирамид, цилиндров и конусов), аппроксимированных многогранниками. Объекты расставлены на ландшафте и служат окружением; пользователь их не изменяет.
	\item Источник света на бесконечности с направлением
	\begin{equation}
		\vec{L} = (\cos\theta_L \cos\varphi_L,\ \cos\theta_L \sin\varphi_L,\ \sin\theta_L),
		\label{eq:light-dir}
	\end{equation}
	где $\varphi_L$ "--- азимут, $\theta_L$ "--- высота над горизонтом.
	\item Не более восьми занавесов северного сияния. Занавесом северного сияния (далее "--- занавесом) в работе называется вертикальный светящийся слой, ограниченный по высоте границами свечения $h_{\min}$ и $h_{\max}$, над основанием "--- кривой на горизонтальной плоскости, вытянутой вдоль основного направления и изгибающейся в виде складок. Протяжённость $\ell$ "--- длина основания вдоль основного направления, амплитуда складок $A$ "--- наибольшее отклонение основания от него, $\lambda$ "--- характерная длина складки; цвета $C_{\text{н}}$, $C_{\text{с}}$, $C_{\text{в}}$ задаются для нижней, средней и верхней третей высоты свечения. Положение и направление занавесов порождаются генератором псевдослучайных чисел.
	\item Наблюдатель с постоянным углом обзора по вертикали $60^\circ$.
\end{enumerate}

Сцена освещается источником и свечением занавесов, тени строятся от обоих. Расчётная сетка покрывает область $[0;\,100] \times [0;\,100] \times [10;\,60]$~у.е., в которой могут находиться занавесы. Карта высот, состав, цвета и расстановка объектов, а также отражающие свойства поверхностей ландшафта и объектов "--- коэффициенты диффузного $k_d$ и зеркального $k_s$ отражения и показатель блеска $n$ "--- загружаются из файлов. Остальные параметры задаются во время работы программы (таблица~\ref{tab:input-ranges}), отображение ландшафта и объектов можно отключить.

\begin{table}[H]
	\caption{Допустимые диапазоны входных параметров}
	\label{tab:input-ranges}
	\begin{tabular}{|L{0.42}|C{0.22}|C{0.35}|}
		\hline
		Параметр & Обозначение & Диапазон \\
		\hline
		азимут и высота источника & $\varphi_L$; $\theta_L$ & $[0^\circ;\,360^\circ)$; $[0^\circ;\,90^\circ]$ \\
		\hline
		интенсивность и цвет источника & $I_L$; $C_L$ & $[0;\,1]$; RGB, $[0;\,1]^3$ \\
		\hline
		цвета частей занавеса & $C_{\text{н}}$, $C_{\text{с}}$, $C_{\text{в}}$ & RGB, $[0;\,1]^3$ \\
		\hline
		количество занавесов & $K$ & $[1;\,8]$ \\
		\hline
		границы свечения & $h_{\min}$, $h_{\max}$ & $10 \leqslant h_{\min} < h_{\max} \leqslant 60$~у.е. \\
		\hline
		протяжённость, амплитуда и длина складки & $\ell$; $A$; $\lambda$ & $[5;\,100]$; $[0;\,20]$; $[1;\,\ell]$~у.е. \\
		\hline
		коэффициенты излучения и поглощения среды & $\varepsilon$; $\kappa$ & $[0;\,10]$; $[0;\,1]$~у.е.$^{-1}$ \\
		\hline
		начальное значение генератора псевдослучайных чисел & $r_0$ & $[0;\,2^{32} - 1]$ \\
		\hline
		положение наблюдателя & $P_V$ & внутри сцены, $z_V > H(x_V, y_V)$ \\
		\hline
		азимут и угол места направления взгляда & $\varphi_V$; $\theta_V$ & $[0^\circ;\,360^\circ)$; $[-90^\circ;\,90^\circ]$ \\
		\hline
		размер расчётной сетки & $N_x$, $N_y$; $N_z$ & $[16;\,256]$; $[8;\,128]$ \\
		\hline
		шаг интегрирования вдоль луча & $\Delta s$ & $[0{,}05;\,5]$~у.е. \\
		\hline
	\end{tabular}
\end{table}

Условие $h_{\min} \geqslant 10$~у.е. исключает пересечение занавесов с рельефом. Коэффициенты отражения в файлах удовлетворяют условию $k_d + k_s \leqslant 1$, так как поверхность не может отразить больше света, чем на неё падает~\cite[с.~385]{rogers1989}. По теореме отсчётов (теореме Котельникова) периодическая составляющая, дискретизованная с шагом $\delta$, передаётся без искажений, только если её период больше $2\delta$~\cite[с.~261--264]{gonzalez2012}; на сетке наибольшего размера $\delta \approx 0{,}39$~у.е., поэтому принято $\lambda \geqslant 1$~у.е.

% Формализация поставленной задачи в нотации IDEF0 представлена на рисунке~\ref{fig:idef0-a0}.

% \begin{figure}[H]
% 	\centering
% 	\IfFileExists{img/idef0-a0.pdf}{%
% 		\includegraphics[width=\linewidth]{img/idef0-a0.pdf}%
% 	}{%
% 		\fbox{\parbox[c][6cm][c]{0.9\linewidth}{\centering Диаграмма IDEF0 A-0 (экспорт из Ramus: img/idef0-a0.pdf)}}%
% 	}
% 	\caption{Функциональная модель программы в нотации IDEF0}
% 	\label{fig:idef0-a0}
% \end{figure}